%% kapa: exam=1 round=23 year=2012 period=2 form=A subject=appraisal_law subject_ko=부동산관계법규 seq=33 number=73
%% source: Q-Net 감정평가사 기출문제 (cst00309) artlSeq=5052580 file=raw/1008/5052580/2012년 제23회 감정평가사 제1차시험 기출문제(2교시 A형).pdf
%% conversion: kordoc 4.24.51 -> tools/gen1.py (auto); review=proofread-sonnet; flags=none
\begin{question}{73}
측량ㆍ수로조사 및 지적에 관한 법령상 지번의 구성 및 부여에 관한 설명으로 옳은 것은?
\begin{choices}
  \item 경계점좌표등록부에는 지번을 등록하지 아니 한다.
  \item 여러 필지로 되어 있는 토지를 등록전환할 경우에는 그 지번부여지역의 최종 본번의 다음 순번부터 본번으로 하여 순차적으로 지번을 부여할 수 있다.
  \item 도시개발사업 시행지역의 지번은 그 사업의 준공 이전에 부여하여야 한다.
  \item 지번은 북동에서 남서로 순차적으로 부여한다.
  \item 지적소관청이 지번부여지역 전부의 지번을 변경할 경우에는 국토해양부장관의 승인을 받고, 일부의 지번을 변경할 경우에는 시ㆍ도지사나 대도시 시장의 승인을 받아야 한다.
\end{choices}
\end{question}
\answer{2}
